\documentclass[14pt,a4paper]{extarticle}

\usepackage[T2A]{fontenc}
\usepackage[utf8]{inputenc}
\usepackage[russian]{babel}
\usepackage[a4paper,margin=2cm]{geometry}
\usepackage{amsmath,amssymb}
\usepackage{graphicx}
\usepackage{booktabs}
\usepackage{longtable}
\usepackage{float}
\usepackage{array}
\usepackage{caption}
\usepackage{hyperref}
\usepackage{setspace}

\onehalfspacing
\graphicspath{{../outputs/lab2/plots/}}

\hypersetup{
    colorlinks=true,
    linkcolor=black,
    urlcolor=blue
}

\newcommand{\plotplaceholder}[2]{
    \begin{figure}[H]
        \centering
        \IfFileExists{../outputs/lab2/plots/#1}{
            \includegraphics[width=0.88\textwidth]{#1}
        }{
            \fbox{
                \parbox[c][5cm][c]{0.82\textwidth}{
                    \centering
                    Файл рисунка \texttt{\detokenize{#1}} создается запуском\\
                    \texttt{python lab2/run.py}
                }
            }
        }
        \caption{#2}
    \end{figure}
}

\title{Лабораторная работа №2\\Классические методы градиентного спуска}
\author{
    Новиков Алексей Георгиевич, группа М3233\\
    Душкин Александр Алексеевич, группа М3233
}
\date{2026}

\begin{document}

\begin{titlepage}
    \begin{center}
        {\large Отчет по лабораторной работе №2}\\[0.5cm]
        {\LARGE \textbf{Градиентный спуск: классика}}\\[2cm]

        \begin{flushright}
            Выполнили:\\[0.2cm]
            Новиков Алексей Георгиевич, группа М3233\\
            Душкин Александр Алексеевич, группа М3233\\[0.5cm]
        \end{flushright}

        \vfill
        Санкт-Петербург\\[0.2cm]
        2026
    \end{center}
\end{titlepage}

\tableofcontents
\newpage

\section{Постановка задачи}

Цель работы --- реализовать и сравнить классические методы градиентной оптимизации для гладких функций двух переменных.
В работе исследуются:
\begin{enumerate}
    \item градиентный спуск с постоянным шагом;
    \item градиентный спуск с дроблением шага по условию Армихо;
    \item градиентный спуск с выбором шага по сильным условиям Вольфе;
    \item метод наискорейшего градиентного спуска с одномерным поиском.
\end{enumerate}

Для каждого метода фиксировались число итераций, количество вычислений функции, количество вычислений градиента, итоговая точка, значение целевой функции и факт выполнения критерия остановки.
Основной критерий остановки задавался нормой градиента:
\[
    \|\nabla f(x_k)\| \le \varepsilon.
\]

\textbf{Промежуточный вывод.} В отличие от методов одномерного поиска из первой лабораторной, здесь каждый шаг состоит из двух логических частей: выбора направления и выбора длины шага. Направление во всех классических вариантах строится по градиенту, а вот длина шага может быть фиксированной, подобранной обратным ходом или найденной одномерной минимизацией. Поэтому качество метода нельзя оценивать только числом итераций: методы с более умным выбором шага обычно делают меньше внешних шагов, но платят за это дополнительными обращениями к функции и градиенту.

\section{Используемые модели задач}

\subsection{Хорошо обусловленная квадратичная функция}

Первая тестовая функция:
\[
    f(x,y)=\frac{1}{2}(x^2+y^2).
\]
Минимум достигается в точке $(0,0)$, значение минимума равно нулю. Матрица Гессе равна единичной, поэтому число обусловленности равно $1$.

\textbf{Промежуточный вывод.} Эта функция используется как базовый эталон: все координатные направления имеют одинаковую кривизну, а антиградиент из любой точки направлен прямо к минимуму. Из-за такой структуры различия между методами выбора шага проявляются почти изолированно: если шаг подобран корректно, траектория не испытывает зигзагов и быстро идет к центру линий уровня.

\subsection{Плохо обусловленная квадратичная функция}

Вторая квадратичная функция:
\[
    g(x,y)=\frac{1}{2}x^2+50y^2.
\]
Ее градиент имеет вид
\[
    \nabla g(x,y) = (x,100y).
\]
Число обусловленности матрицы Гессе равно $100$.

\textbf{Промежуточный вывод.} Плохая обусловленность приводит к вытянутым эллиптическим линиям уровня, потому что по координате $y$ кривизна в сто раз больше, чем по координате $x$. Градиент указывает направление наискорейшего локального убывания, но не учитывает, что вдоль разных осей допустимы разные масштабы движения. Поэтому слишком большой общий шаг легко перескакивает через узкую долину, а слишком малый делает продвижение вдоль пологой координаты медленным.

\subsection{Сложные функции}

Для проверки поведения методов за пределами простых выпуклых квадратичных моделей использовались функции Розенброка, Экли и Химмельблау.

\paragraph{Функция Розенброка.}
\[
    r(x,y)=(1-x)^2+100(y-x^2)^2.
\]
Глобальный минимум достигается в точке $(1,1)$.

\paragraph{Функция Экли.}
\[
    a(x,y)=-20\exp\left(-0.2\sqrt{\frac{x^2+y^2}{2}}\right)
    -\exp\left(\frac{\cos(2\pi x)+\cos(2\pi y)}{2}\right)+e+20.
\]
Глобальный минимум находится в точке $(0,0)$, однако поверхность содержит большое количество локальных особенностей.

\paragraph{Функция Химмельблау.}
\[
    h(x,y)=(x^2+y-11)^2+(x+y^2-7)^2.
\]
Функция имеет четыре глобальных минимума.

\textbf{Промежуточный вывод.} Эти функции проверяют разные трудности, напрямую связанные с логикой градиентных методов. У Розенброка антиградиент часто направлен поперек узкой долины, поэтому метод вынужден делать много корректирующих шагов. У Экли локальные колебания могут менять направление градиента и уводить метод в соседние области. У Химмельблау несколько минимумов, поэтому результат зависит не только от метода, но и от того, в какую область притяжения попадает стартовая точка.

\section{Описание реализованных методов}

\subsection{Градиентный спуск с постоянным шагом}

Метод использует обновление
\[
    x_{k+1}=x_k-\alpha\nabla f(x_k),
\]
где $\alpha$ фиксируется заранее. В лабораторной работе для квадратичных функций перебирались значения
\[
    \alpha \in \{10^{-4},3\cdot10^{-4},10^{-3},3\cdot10^{-3},10^{-2},10^{-1}\}.
\]

\textbf{Промежуточный вывод.} Постоянный шаг прост и дешев на итерацию, потому что после вычисления градиента метод сразу делает обновление и не проверяет качество выбранного $\alpha$. Именно отсутствие обратной связи делает его чувствительным к масштабу функции: слишком малый шаг дает почти правильное, но очень медленное движение, а слишком большой шаг не проверяется условием убывания и может привести к расходимости, особенно при большой кривизне по одной из координат.

\subsection{Дробление шага по условию Армихо}

На каждой итерации выбирается направление $p_k=-\nabla f(x_k)$, после чего начальный шаг $\alpha_0$ последовательно уменьшается:
\[
    \alpha \leftarrow \rho\alpha, \qquad 0<\rho<1,
\]
пока не выполнится условие достаточного убывания:
\[
    f(x_k+\alpha p_k)\le f(x_k)+c_1\alpha\nabla f(x_k)^Tp_k.
\]

\textbf{Промежуточный вывод.} Условие Армихо защищает метод от слишком длинных шагов за счет явной проверки достаточного убывания функции. Логика метода устроена консервативно: если пробный шаг не дает нужного уменьшения, он дробится до приемлемого масштаба. Поэтому алгоритм устойчивее фиксированного шага, но каждое неудачное дробление означает новое вычисление функции.

\subsection{Сильные условия Вольфе}

Сильные условия Вольфе требуют одновременно достаточного убывания и ограничения производной вдоль направления:
\[
    f(x_k+\alpha p_k)\le f(x_k)+c_1\alpha\nabla f(x_k)^Tp_k,
\]
\[
    |\nabla f(x_k+\alpha p_k)^Tp_k|\le c_2|\nabla f(x_k)^Tp_k|.
\]

В реализации используется фаза расширения интервала и процедура \texttt{zoom}, близкая к алгоритму из Носедала--Райта.

\textbf{Промежуточный вывод.} Сильные условия Вольфе выбирают шаг более осмысленно, чем одно только условие Армихо, потому что контролируют не только уменьшение функции, но и наклон вдоль выбранного направления. Это не дает принять слишком короткий шаг, при котором функция формально уменьшилась, но движение почти остановилось. Такая проверка требует вычислять производную вдоль направления, поэтому метод расходует больше обращений к градиенту.

\subsection{Наискорейший градиентный спуск}

В методе наискорейшего спуска направление остается антиградиентным, но длина шага находится решением одномерной задачи:
\[
    \alpha_k=\arg\min_{\alpha\ge0} f(x_k-\alpha\nabla f(x_k)).
\]
В библиотеке для этой одномерной подзадачи используется метод золотого сечения.

\textbf{Промежуточный вывод.} Наискорейший спуск часто резко уменьшает число внешних итераций, потому что для текущего антиградиентного направления выбирает почти оптимальную длину шага, а не просто приемлемую. Однако структура метода переносит сложность внутрь итерации: одномерный поиск многократно вычисляет целевую функцию вдоль прямой, поэтому выигрыш по числу шагов не всегда означает минимальную вычислительную стоимость.

\section{Структура программной реализации}

Лабораторная работа использует общую библиотеку \texttt{optimlib}. Основные элементы реализации:
\begin{itemize}
    \item \texttt{MultivariateFunction} --- базовый класс многомерной функции с подсчетом вызовов функции и градиента;
    \item \texttt{GradientOptimizer} --- общий каркас классических градиентных методов;
    \item \texttt{ConstantStepGD}, \texttt{ArmijoBacktracking}, \texttt{StrongWolfe}, \texttt{SteepestDescent} --- конкретные оптимизаторы;
    \item \texttt{OptimizationExperiment} --- запуск сетки экспериментов по функциям, методам, точностям и параметрам;
    \item \texttt{HistoryCallback} --- накопление траектории оптимизации для построения линий уровня и графиков сходимости.
\end{itemize}

Конфигурация экспериментов вынесена в \texttt{lab2/config.yaml}. Это позволяет изменять набор функций, точностей и параметров без изменения кода оптимизаторов.

\textbf{Промежуточный вывод.} Разделение на библиотеку и конфигурационные файлы соответствует структуре самих экспериментов: математическая логика метода находится в оптимизаторе, параметры запуска --- в конфиге, а сбор статистики --- в общем раннере. Благодаря этому один и тот же метод можно запускать на разных функциях и точностях без дублирования кода, а сравнение остается честным, потому что все результаты проходят через единый формат.

\section{Результаты вычислительных экспериментов}

\subsection{Подбор постоянного шага на квадратичных функциях}

В таблице приведены результаты для точности $\varepsilon=10^{-8}$.

\begin{table}[H]
\centering
\caption{Градиентный спуск с постоянным шагом на квадратичных функциях}
\begin{tabular}{llrrrr}
\toprule
Функция & $\alpha$ & Сошелся & Итерации & $f(x_k)$ & Вызовы градиента\\
\midrule
$f$ & $10^{-4}$ & нет & 2000 & $2.681$ & 2001\\
$f$ & $3\cdot10^{-4}$ & нет & 2000 & $1.205$ & 2001\\
$f$ & $10^{-3}$ & нет & 2000 & $7.31\cdot10^{-2}$ & 2001\\
$f$ & $3\cdot10^{-3}$ & нет & 2000 & $2.41\cdot10^{-5}$ & 2001\\
$f$ & $10^{-2}$ & да & 1937 & $4.93\cdot10^{-17}$ & 1938\\
$f$ & $10^{-1}$ & да & 185 & $4.70\cdot10^{-17}$ & 186\\
$g$ & $10^{-4}$ & нет & 2000 & $1.341$ & 2001\\
$g$ & $3\cdot10^{-4}$ & нет & 2000 & $6.02\cdot10^{-1}$ & 2001\\
$g$ & $10^{-3}$ & нет & 2000 & $3.66\cdot10^{-2}$ & 2001\\
$g$ & $3\cdot10^{-3}$ & нет & 2000 & $1.21\cdot10^{-5}$ & 2001\\
$g$ & $10^{-2}$ & да & 1902 & $4.98\cdot10^{-17}$ & 1903\\
$g$ & $10^{-1}$ & ошибка & -- & -- & --\\
\bottomrule
\end{tabular}
\end{table}

\textbf{Промежуточный вывод.} Для хорошо обусловленной функции лучший из протестированных шагов --- $\alpha=0.1$, потому что одинаковая кривизна по координатам позволяет делать достаточно длинное движение вдоль антиградиента. Для плохо обусловленной функции такой шаг становится слишком большим: компонента градиента по $y$ масштабируется множителем $100$, поэтому обновление перескакивает через минимум в крутом направлении. Устойчивым оказывается меньший шаг $\alpha=0.01$, что согласуется с тем, что допустимый постоянный шаг ограничен наибольшим собственным значением матрицы Гессе.

\subsection{Методы выбора шага на квадратичных функциях}

\begin{table}[H]
\centering
\caption{Сравнение методов при $\varepsilon=10^{-8}$ на квадратичных функциях}
\begin{tabular}{llrrrr}
\toprule
Функция & Метод & Итерации & Вызовы $f$ & Вызовы $\nabla f$ & $f(x_k)$\\
\midrule
$f$ & Armijo & 1 & 3 & 2 & $0$\\
$f$ & Strong Wolfe & 1 & 3 & 3 & $0$\\
$f$ & Steepest Descent & 1 & 46 & 2 & $1.11\cdot10^{-17}$\\
$g$ & Armijo & 869 & 6595 & 870 & $2.77\cdot10^{-17}$\\
$g$ & Strong Wolfe & 869 & 7464 & 1783 & $2.77\cdot10^{-17}$\\
$g$ & Steepest Descent & 11 & 484 & 12 & $1.47\cdot10^{-20}$\\
\bottomrule
\end{tabular}
\end{table}

\textbf{Промежуточный вывод.} На хорошо обусловленной функции все методы быстро попадают в минимум, потому что направление антиградиента хорошо согласовано с направлением на минимум, а выбор шага почти не сталкивается с конфликтом масштабов. На плохо обусловленной функции наискорейший спуск выигрывает по числу внешних итераций: он каждый раз подбирает шаг, минимизирующий функцию вдоль текущего направления, и поэтому лучше проходит вытянутую долину. Но этот выигрыш достигается внутренним одномерным поиском, из-за чего число вычислений функции на одну внешнюю итерацию значительно выше.

\subsection{Сложные функции}

В таблице показано лучшее значение функции среди запусков из заданных стартовых точек при $\varepsilon=10^{-8}$.

\begin{table}[H]
\centering
\caption{Результаты на сложных функциях}
\begin{tabular}{llrrrr}
\toprule
Функция & Метод & Успешных запусков & Мин. итераций & Мин. вызовов $f$ & Лучшее $f(x_k)$\\
\midrule
Ackley & ConstantStepGD & 12 & 27 & 28 & $2.95$\\
Ackley & Armijo & 0 & 2000 & 81669 & $1.71\cdot10^{-12}$\\
Ackley & Strong Wolfe & 0 & 2000 & 85963 & $3.22\cdot10^{-12}$\\
Ackley & Steepest Descent & 0 & 2000 & 86001 & $1.64\cdot10^{-8}$\\
Rosenbrock & ConstantStepGD & 0 & 2000 & 2001 & $1.39\cdot10^{-2}$\\
Rosenbrock & Armijo & 0 & 2000 & 21345 & $1.29\cdot10^{-4}$\\
Rosenbrock & Strong Wolfe & 0 & 2000 & 23348 & $1.31\cdot10^{-6}$\\
Rosenbrock & Steepest Descent & 0 & 2000 & 86001 & $5.82\cdot10^{-6}$\\
Himmelblau & ConstantStepGD & 10 & 22 & 23 & $1.90\cdot10^{-19}$\\
Himmelblau & Armijo & 3 & 17 & 141 & $4.26\cdot10^{-20}$\\
Himmelblau & Strong Wolfe & 3 & 17 & 158 & $6.75\cdot10^{-20}$\\
Himmelblau & Steepest Descent & 3 & 7 & 302 & $8.70\cdot10^{-21}$\\
\bottomrule
\end{tabular}
\end{table}

\textbf{Промежуточный вывод.} На функции Розенброка все классические методы испытывают трудности, потому что локальное направление антиградиента часто указывает поперек узкой долины, а не вдоль нее. Методы выбора шага могут подобрать безопасную длину движения, но они не меняют само направление на более криволинейно согласованное, поэтому сходимость остается медленной. На функции Химмельблау методы работают лучше: в выбранных стартовых точках градиентные направления приводят в области притяжения глобальных минимумов, а поверхность менее похожа на длинный узкий канал.

\subsection{Графики траекторий и сходимости}

\plotplaceholder{WellConditionedQuadratic_trajectories.png}{Траектории методов на хорошо обусловленной квадратичной функции}
\plotplaceholder{IllConditionedQuadratic_trajectories.png}{Траектории методов на плохо обусловленной квадратичной функции}
\plotplaceholder{Rosenbrock__x0_-1.2_1.0__trajectories.png}{Траектории методов на функции Розенброка}
\plotplaceholder{Himmelblau__x0_2.0_2.0__trajectories.png}{Траектории методов на функции Химмельблау}

\textbf{Промежуточный вывод.} Траектории на линиях уровня позволяют увидеть, почему одинаковое число итераций не является универсальной мерой качества. Если метод использует только антиградиент и фиксированный шаг, он может многократно пересекать узкую долину короткими зигзагообразными движениями. Если же шаг выбирается линейным поиском, траектория становится более осмысленной относительно текущего направления, но часть работы скрыта внутри процедуры подбора шага и проявляется в увеличенном числе вызовов функции.

\section{Сводная таблица эксперимента}

Полная сводка всех запусков сохраняется в файле \texttt{../outputs/lab2/summary.csv}. Если доступна LaTeX-версия таблицы, она подключается ниже.

\IfFileExists{../outputs/lab2/summary.tex}{
    \scriptsize
    \input{../outputs/lab2/summary.tex}
    \normalsize
}{
    \begin{center}
        \fbox{\parbox{0.82\textwidth}{Полная таблица доступна в \texttt{../outputs/lab2/summary.csv}.}}
    \end{center}
}

\section{Общий вывод}

В ходе лабораторной работы была реализована и исследована группа классических градиентных методов. Эксперименты подтвердили основные теоретические свойства этих алгоритмов.

Градиентный спуск с постоянным шагом оказался самым дешевым по одной итерации, но самым чувствительным к параметру $\alpha$. На хорошо обусловленной квадратичной функции большой шаг $\alpha=0.1$ дал быструю сходимость, тогда как на плохо обусловленной функции тот же шаг привел к ошибке вычислений.

Методы дробления шага по Армихо и сильным условиям Вольфе показали большую устойчивость. Они автоматически подбирают допустимую длину шага, однако делают это за счет дополнительных вычислений функции и, в случае Вольфе, градиента.

Метод наискорейшего спуска оказался особенно эффективным на плохо обусловленной квадратичной функции по числу внешних итераций: $11$ итераций против $869$ у методов Армихо и Вольфе. Однако его стоимость переносится во внутренний одномерный поиск, поэтому число вычислений функции остается заметным.

На сложных функциях различия стали более выраженными. Функция Розенброка оказалась наиболее трудной из-за узкой изогнутой долины; функция Химмельблау, напротив, хорошо решалась из выбранных стартовых точек. Таким образом, универсального лучшего метода среди рассмотренных нет: выбор алгоритма зависит от геометрии функции и от допустимой стоимости выбора шага.

\end{document}
